The measurement in Section~\ref{sec:results} needs an explanation. Under a
correctly specified model the inverse-Hessian standard error is consistent for the
sampling standard deviation, so a pooled ratio of $q = \HsixQ{}$ should not appear
as an asymptote. This section derives the object the library actually reports,
states what has to break for it to understate its own spread, and sets a
leading-order correction against the measured per-block ratios. The correction
does not reproduce the measurement, and the size of the disagreement is itself the
most informative number in this section.

\subsection{The asymptotic covariance of the estimator}

Write $\ell_T(\theta)$ for \eqref{eq:likelihood} and $U_T(\theta) = \nabla\ell_T(\theta)$
for the score, with $p = M + 2M^2$ coordinates. Under stationarity and the usual
regularity conditions for point-process likelihoods, the maximum-likelihood
estimator satisfies
\begin{equation}
\sqrt{N}\,\big(\hat\theta - \theta\big) \;\overset{d}{\longrightarrow}\;
\mathcal{N}\!\left(0,\; I(\theta)^{-1} J(\theta) I(\theta)^{-1}\right),
\label{eq:sandwich}
\end{equation}
where $I(\theta) = \mathbb{E}[-\nabla^2 \ell_T(\theta)]$ is the expected
information per replication and $J(\theta) = \mathbb{E}[U_T(\theta) U_T(\theta)^{\top}]$
is the covariance of the score. The matrix in \eqref{eq:sandwich} is the sandwich,
and the library reports its middle factor inverted: the observed information
$-\nabla^2 \ell_T(\hat\theta)$, evaluated at the estimate rather than averaged
over the truth.

The gap between the two is governed by whether $I = J$. For a point process whose
intensity is the true conditional intensity, the score is a martingale with
respect to the natural filtration, and the predictable-variation identity gives
\begin{equation}
J(\theta) = \mathbb{E}\!\left[\int_0^T \sum_{i} \lambda_i(u)\,
\nabla\log\lambda_i(u)\,\nabla\log\lambda_i(u)^{\top} du\right] = I(\theta)
\label{eq:equality}
\end{equation}
for every $\theta$ in the interior of the parameter space. Equation~\eqref{eq:equality}
is the information equality in its point-process form: under exact specification the
observed information is an unbiased estimator of the score variance, and the
sandwich collapses to $I^{-1}$, which is what the library computes. When the model
is correct and the window is long against the decay time, then, the asymptotic
theory predicts $q \to 1$. The shortfall measured here is not a misspecified
intensity in the usual sense; it comes from the finite window, from parameters whose
information the window cannot supply, and from the arithmetic of comparing a
reported interval with a realised spread.

\subsection{What breaks the equality}

Three violations of \eqref{eq:equality} appear at finite window, and a fourth effect
in the ratio itself is larger than all three. The first three are worth separating
because they are shared with ordinary finite-sample theory; the fourth is specific
to measuring a reported interval against a realised spread.

The likelihood is not the sampling distribution of a fixed number of events. The
study fixes the window $T$ and lets the event count fall where it may, and a
stationary Hawkes process observed on $[0, T]$ has a random count whose mean is
$T$ times the mean intensity. The information in \eqref{eq:equality} therefore
scales with $T$ rather than being held fixed by the design, and the realised count
is itself informative about $\rho$, so a reading that treats the count as ancillary
is reading an approximation. The study reports its realised counts as within
\ThCountTolerance{} percent of target, which bounds this effect without removing
it. It is present for every parameter, and it is the standard reason a likelihood
conditioned on an achieved count understates the variance of its own estimate.

The exponential kernel truncates information at the window edge. An event at time
$t$ contributes to the intensity at $u > t$ with weight $e^{-\beta(u - t)}$, so the
fraction of its kernel mass that falls outside the window is $1 - e^{-\beta(T-t)}$.
For an event in the last $1/\beta$ of the window that fraction is at least
$1 - e^{-1} \approx 0.63$ and can be arbitrarily close to one at the boundary,
while for an event early in the window it is negligible. The effective sample size
for the decay parameters is therefore closer to $\beta_{\min} T$ than to $N$, and
the two are very different across this grid: $\beta_{\min} T$ is between
\ThBetaMinTLow{} and roughly \ThBetaMinTHigh{} in the exact-timestamp cells at
$c = 1$, and below one in \ThShallowBelowOne{} of the \ThShallowCells{}
exact-timestamp shallow-decay cells.

The decay parameters are weakly identified when $\beta_{\min} T$ is small. The
information about $\beta_{ij}$ comes from the curvature of the decay profile
within the window, and if the window holds less than one decay time, the profile
is nearly flat and the information nearly vanishes. In the limit the information
matrix is singular in the $\beta$ direction, and any inverse is unstable.

The reported interval is a nonlinear functional of the information. This is the
one that matters, and it is the subject of the next subsection.

\subsection{The leading-order correction}

The quantity the study measures is a ratio of averages, $q_B = \text{se}_B /
\text{sd}_B$, and the two averages are taken over the same replications but not
over the same functional. The numerator averages the reported standard error,
which is $\sqrt{\widehat{\text{Var}}}$ evaluated at that replication's estimate;
the denominator takes the square root of an averaged squared deviation. For a
coordinate $j$, write $X_r = \widehat{\text{Var}}_j(\hat\theta_r)$ for the reported
variance at replication $r$. The reported standard error is $\sqrt{X_r}$ and the
squared spread the denominator measures is, to leading order,
$\mathbb{E}[X_r]$ plus a term of order $N^{-1}$ that the bias expansion supplies.
Then
\begin{equation}
q_j \;=\; \frac{\mathbb{E}[\sqrt{X_j}]}{\sqrt{\mathbb{E}[X_j]} + O(N^{-1})}.
\label{eq:qjensen}
\end{equation}
The average of a concave function is at most the function of the average, so
$\mathbb{E}[\sqrt{X_j}] \le \sqrt{\mathbb{E}[X_j]}$, with equality only when $X_j$
is constant across replications. Because $\sqrt{\cdot}$ is strictly concave, the
inequality is strict whenever the reported variance varies at all, and the size of
the gap is set by that variation. Expand $\sqrt{X}$ about $\mathbb{E}[X]$ with
$X = \mu_X(1 + \epsilon)$, $\mathbb{E}[\epsilon] = 0$:
\begin{equation}
\frac{\mathbb{E}[\sqrt{X}]}{\sqrt{\mu_X}}
= \mathbb{E}\big[(1 + \epsilon)^{1/2}\big]
= 1 - \tfrac{1}{8}\operatorname{Var}(\epsilon) + O(\mathbb{E}[\epsilon^3]),
\label{eq:jensen}
\end{equation}
so to leading order
\begin{equation}
q_j \;\approx\; \frac{1}{1 + \operatorname{CV}(X_j)^2/8}
\;\approx\; 1 - \frac{\operatorname{CV}(X_j)^2}{8},
\qquad \operatorname{CV}(X_j) = \frac{\operatorname{sd}(X_j)}{\mathbb{E}[X_j]}.
\label{eq:qvarcv}
\end{equation}
The reported interval systematically understates the spread whenever the reported
variance is not constant, and the shortfall is one eighth of its squared
coefficient of variation.

Equation~\eqref{eq:qvarcv} can be inverted at the block level, and inverted it is a
description of the reported variance's dispersion rather than a test of it, since
the right-hand side is fitted from the same numbers as the left. Inverting at the
three measured block ratios gives implied coefficients of variation of
\ThCvMu{} for the baselines, \ThCvAlpha{} for the excitation entries and
\ThCvBeta{} for the decay rates. The decay block's reported variance therefore
varies \ThCvRatio{} times as much as the baseline's. The deficit ratio
\ThDefRatio{} carries no information beyond that one, since \eqref{eq:qvarcv} makes
it the square of the CV ratio by construction. What the inversion does not do is
predict which block should have the larger coefficient of variation; the identity
holds for every $q$, and the prediction that the decay block is the irregular one
comes from the ridge argument below, not from this equation.

The relevant scale for a regular parameter is $\operatorname{CV}(X_j) = O(N^{-1/2})$,
because the information is an average of $N$ terms and its relative dispersion
shrinks at that rate. The Poisson baseline makes the constant concrete: for a
single type with a known decay or a baseline-only model, the estimated rate is the
event count over $T$, its squared coefficient of variation is exactly $1/N$, and
\eqref{eq:qvarcv} gives a deficit of $1/(8N)$ with no free constant. That is the
most generous regularity benchmark this design admits, and at the smallest pooled
event count it predicts a deficit of \ThRegDeficit{}, against a measured decay
deficit of \ThDefBeta{}. The two differ by a factor of \ThMismatch. The Jensen
correction therefore explains the direction and the ordering across blocks and
cannot supply the magnitude in the decay block; the decay block is not behaving
like a regular parameter, because its information is not an average over a large
effective sample but a functional of a nearly flat direction whose curvature is
itself random. The derivation closes to \eqref{eq:qvarcv} and stops there. A
closed-form prediction of the ratio would need the joint distribution of the
reported variance's fluctuation in the decay direction, which this study did not
record and cannot recover from the committed outputs.

\subsection{Why the decay block is not regular}

The failure of the $O(N^{-1/2})$ scale in the decay block is not a matter of
constants. It follows from the shape of the information matrix near the ridge, and
it is worth writing out because it is what makes the whole phenomenon have a heavy
tail rather than a uniform shrinkage.

Write the information matrix in its eigenbasis, $I(\theta) = \sum_k \lambda_k
v_k v_k^{\top}$ with $\lambda_1 \ge \cdots \ge \lambda_p > 0$. The variance of a
coordinate $j$ is the diagonal of the inverse,
\begin{equation}
\operatorname{Var}(\hat\theta_j) = \sum_{k=1}^{p} \frac{v_{kj}^2}{\lambda_k},
\label{eq:eigenvar}
\end{equation}
so a coordinate that loads on a nearly null direction has a variance that no
average over events explains. What can be derived here is the direction the decay
parameters share with their own excitation weights, and it is derived from the
likelihood rather than asserted.

Fix one pair $(i,j)$ and consider the path that sends $\beta_{ij} \to \infty$ while
holding $\alpha_{ij}/\beta_{ij} = n_{ij}$ constant, so $\alpha_{ij} = n_{ij}
\beta_{ij}$ grows with it. The excitation term in the intensity at an event time
$t_k$ is $\sum_{l} \alpha_{ij} e^{-\beta_{ij}(t_k - t_l)} = n_{ij} \beta_{ij}
\sum_l e^{-\beta_{ij}(t_k - t_l)}$, and for exact timestamps no two events coincide,
so for every gap $\Delta > 0$ this tends to zero as $\beta_{ij} \to \infty$. The
compensator integral tends instead to $n_{ij}$ times the number of type-$j$ events,
because the kernel integrates to $\alpha_{ij}/\beta_{ij} = n_{ij}$ however narrow it
becomes. The log-likelihood along this path therefore converges to a limit in which
$\beta_{ij}$ appears only through the fixed ratio $n_{ij}$, so its derivative with
respect to $\beta_{ij}$ vanishes and the direction is flat. Along it the two
parameters move together at constant ratio, not in opposite directions, and the
baselines do not move at all to first order: the stationary mean intensity is
$(I - \Gamma)^{-1}\mu$ and $\Gamma_{ij} = \alpha_{ij}/\beta_{ij} = n_{ij}$ is what
the path holds fixed.

That derivation establishes a flat direction and the shape of it. It does not
establish that this particular direction is the least-curved eigenvector of the
information matrix at the estimates this study produced, and the committed outputs
cannot settle that, because they record no per-replication Hessian or eigenvector.
Throughout this subsection the identification of \eqref{eq:eigenvar}'s small
eigenvalue with the alpha--beta direction is therefore an assumption that the study
is consistent with rather than a result it proves. What the study does show is that
the decay block is the one whose reported ratio fails, which is what an assumption
of that kind predicts, and that at the anchor cell the coordinate with the smallest
reported ratio in both the excitation and the decay block is the same pair,
$(\alpha_{01}, \beta_{01})$, which is also the pair with the smallest true
excitation weight and therefore the smallest kernel mass along the flat path. That
is consistent with the plateau argument and is not a substitute for the eigenvector.

The baseline block is regular for a structural reason: its parameters do not appear
in the flat direction above, so its information is an average over $N$
contributions and its reported variance has the $O(N^{-1/2})$ dispersion the
expansion assumes. The decay block is irregular for a structural reason as well: it
loads on a direction whose curvature is small and, because the plateau is not
reached exactly at any finite $\beta$, variable from one replication to the next.

That variability is the quantity \eqref{eq:qvarcv} needs and this study does not
have. The reported variance at a replication is not $\lambda_{\min}^{-1}$ evaluated
at the truth, which is a fixed number; it is the curvature at the estimate the
optimizer found, and the estimate lands at a random point along a direction whose
curvature changes as it moves. The average of the reported standard error is
dominated by the many replications where the curvature is typical, and the square
root of the averaged squared deviation is dominated by the few where the estimate
has run furthest along the flat direction. That asymmetry is the reason the decay
block's deficit is so much larger than the correction allows for. It is also the
reason the median-based companion does not single the decay block out: that
companion replaces the mean reported standard error with a median and the standard
deviation with a scaled median absolute deviation, so it is driven by the bulk of
the replications rather than the tail, and across the twelve pooled cells it stays
between \ThQrobBetaLow{} and \ThQrobBetaHigh{} in the decay block, against
\ThQrobMuLow{} to \ThQrobMuHigh{} in the baseline block. Close to one everywhere is
the honest reading. The pooled mean-based ratio is a decay result because the decay
block is where the mean and the median diverge, not because the decay block is
badly calibrated by every measure.

None of this makes the sandwich inconsistent. Equation~\eqref{eq:sandwich} remains
the right target in the limit, and a sandwich estimator computed from a correctly
evaluated $J$ is the natural candidate for closing the gap measured here. The reason
this paper does not report that experiment is that it requires per-replication
score variances, which the study did not store and which cannot be reconstructed
from the committed analysis outputs.

The pooled figure is consistent with the block account and adds nothing to it. The
pooled statistic is the coordinate-weighted mean of the three block ratios, so it
equals $(2 q_\mu + 4 q_\alpha + 4 q_\beta)/10$ exactly, and the pooled deficit of
\ThPooledDeficit{} decomposes into \ThShareMu{}~percent from the baselines,
\ThShareAlpha{}~percent from the excitation entries and \ThShareBeta{}~percent from
the decay rates. The headline is a decay-rate result with a baseline block along
for the ride, and reporting it pooled is a decision about presentation rather than
about evidence.
